%%
%% This is file `sample-sigconf.tex',
%% generated with the docstrip utility.
%%
%% The original source files were:
%%
%% samples.dtx  (with options: `all,proceedings,bibtex,sigconf')
%% 
%% IMPORTANT NOTICE:
%% 
%% For the copyright see the source file.
%% 
%% Any modified versions of this file must be renamed
%% with new filenames distinct from sample-sigconf.tex.
%% 
%% For distribution of the original source see the terms
%% for copying and modification in the file samples.dtx.
%% 
%% This generated file may be distributed as long as the
%% original source files, as listed above, are part of the
%% same distribution. (The sources need not necessarily be
%% in the same archive or directory.)
%%
%%
%% Commands for TeXCount
%TC:macro \cite [option:text,text]
%TC:macro \citep [option:text,text]
%TC:macro \citet [option:text,text]
%TC:envir table 0 1
%TC:envir table* 0 1
%TC:envir tabular [ignore] word
%TC:envir displaymath 0 word
%TC:envir math 0 word
%TC:envir comment 0 0
%%
%% The first command in your LaTeX source must be the \documentclass
%% command.
%%
%% For submission and review of your manuscript please change the
%% command to \documentclass[manuscript, screen, review]{acmart}.
%%
%% When submitting camera ready or to TAPS, please change the command
%% to \documentclass[sigconf]{acmart} or whichever template is required
%% for your publication.
%%
%%
\documentclass[sigconf]{acmart}
\usepackage{float}
%%
%% \BibTeX command to typeset BibTeX logo in the docs
\AtBeginDocument{%
  \providecommand\BibTeX{{%
Bib\TeX}}}

\acmConference[GECCO '26]{Genetic and Evolutionary Computation Conference Companion}{July 13--17, 2026}{San Jose, Costa Rica}
\acmYear{2026}
\copyrightyear{2026}
\acmDOI{10.1145/nnnnnnn.nnnnnnn}
\acmISBN{978-x-xxxx-xxxx-x/YY/MM}

% \setcopyright{acmlicensed}
% \copyrightyear{2018}
% \acmYear{2018}
% \acmDOI{XXXXXXX.XXXXXXX}
% %% These commands are for a PROCEEDINGS abstract or paper.
% \acmConference[Conference acronym 'XX]{Make sure to enter the correct
%   conference title from your rights confirmation email}{June 03--05,
%   2018}{Woodstock, NY}
% %%
% %%  Uncomment \acmBooktitle if the title of the proceedings is different
% %%  from ``Proceedings of ...''!
% %%
% %%\acmBooktitle{Woodstock '18: ACM Symposium on Neural Gaze Detection,
% %%  June 03--05, 2018, Woodstock, NY}
% \acmISBN{978-1-4503-XXXX-X/2018/06}


%%
%% Submission ID.
%% Use this when submitting an article to a sponsored event. You'll
%% receive a unique submission ID from the organizers
%% of the event, and this ID should be used as the parameter to this command.
%%\acmSubmissionID{123-A56-BU3}

%%
%% For managing citations, it is recommended to use bibliography
%% files in BibTeX format.
%%
%% You can then either use BibTeX with the ACM-Reference-Format style,
%% or BibLaTeX with the acmnumeric or acmauthoryear sytles, that include
%% support for advanced citation of software artefact from the
%% biblatex-software package, also separately available on CTAN.
%%
%% Look at the sample-*-biblatex.tex files for templates showcasing
%% the biblatex styles.
%%

%%
%% The majority of ACM publications use numbered citations and
%% references.  The command \citestyle{authoryear} switches to the
%% "author year" style.
%%
%% If you are preparing content for an event
%% sponsored by ACM SIGGRAPH, you must use the "author year" style of
%% citations and references.
%% Uncommenting
%% the next command will enable that style.
%%\citestyle{acmauthoryear}


%%
%% end of the preamble, start of the body of the document source.
\begin{document}

%%
%% The "title" command has an optional parameter,
%% allowing the author to define a "short title" to be used in page headers.
\title{Audience-Customized Translation of Rule-Based Evidence with Large Language Models Across Multiplexer Benchmarks}


%%
%% The "author" command and its associated commands are used to define
%% the authors and their affiliations.
%% Of note is the shared affiliation of the first two authors, and the
%% "authornote" and "authornotemark" commands
%% used to denote shared contribution to the research.

\author{Harsh Bandhey}
\email{harsh.bandhey@cshs.org}
\orcid{0000-0002-4113-0616}
\affiliation{%
  \institution{Cedars Sinai Health Sciences University}
  \city{Los Angeles}
  \state{CA}
  \country{USA}
}

\author{Gabriel Lipschutz-Villa}
\email{gabelip@sas.upenn.edu}
\orcid{0009-0008-0389-1984}
\affiliation{%
  \institution{University of Pennsylvania}
  \city{Philadelphia}
  \state{PA}
  \country{USA}
}

\author{Khoi Dinh}
\email{khoidinh@seas.upenn.edu}
\orcid{0009-0007-6499-9967}
\affiliation{%
  \institution{University of Pennsylvania}
  \city{Philadelphia}
  \state{PA}
  \country{USA}
}

\author{Malek Kamoun}
\email{malekkam@pennmedicine.upenn.edu}
\orcid{0000-0002-2568-1733}
\affiliation{%
  \institution{University of Pennsylvania}
  \city{Philadelphia}
  \state{PA}
  \country{USA}
}

\author{Ryan J. Urbanowicz}
\authornote{Corresponding Author}
\email{ryan.urbanowicz@cshs.org}
\orcid{0000-0002-0487-5555}
\affiliation{%
  \institution{Cedars Sinai Health Sciences University}
  \city{Los Angeles}
  \state{CA}
  \country{USA}
}

%%
%% By default, the full list of authors will be used in the page
%% headers. Often, this list is too long, and will overlap
%% other information printed in the page headers. This command allows
%% the author to define a more concise list
%% of authors' names for this purpose.
\renewcommand{\shortauthors}{Bandhey et al.}

%%
%% The abstract is a short summary of the work to be presented in the
%% article.
\begin{abstract}
Rule-based machine learning models expose prediction-time evidence as matching IF--THEN rules, but raw rule outputs are often too technical for their intended users. We study a constrained pipeline in which a rule-based model supplies structured evidence and a large language model (LLM) translates that evidence for different audiences rather than generating explanations from scratch. We instantiate the pipeline with HEROS on MUX6, MUX11, and MUX20, compare two HEROS training profiles, two prompt experiments, and three audience targets, and generate 1{,}800 explanations. Across all six dataset/profile runs, \emph{Evidence Precision} and \emph{Prediction--Explanation Agreement} remain perfect. Training profile is a major factor: extensive HEROS training sharply improves MUX11, but not MUX20. The glossary-style prompt experiment usually improves audience-fit scores and sometimes lowers \emph{Causal Overclaim Rate}, but it also lengthens outputs and often raises \emph{Hallucination Rate}. Audience customization creates stable trade-offs: layman outputs minimize hallucination and overclaim, clinician outputs remain the most causally overstated, and expert outputs preserve the most evidence. These preliminary results suggest that constrained LLM translation is feasible, but its success depends jointly on rule-model quality, prompt design, and audience target.

\end{abstract}

%%
%% The code below is generated by the tool at http://dl.acm.org/ccs.cfm.
%% Please copy and paste the code instead of the example below.
%%
\ccsdesc[500]{Computing methodologies~Machine learning}
\ccsdesc[300]{Computing methodologies~Rule learning}
\ccsdesc[300]{Human-centered computing~Natural language interfaces}
\ccsdesc[100]{Computing methodologies~Natural language generation}

%%
%% Keywords. The author(s) should pick words that accurately describe
%% the work being presented. Separate the keywords with commas.
\keywords{interpretable machine learning, rule-based learning, natural language explanations, large language models, explainable AI, evaluation}

%% A "teaser" image appears between the author and affiliation
%% information and the body of the document, and typically spans the
%% page.
% \begin{teaserfigure}
%   \includegraphics[width=\textwidth]{sampleteaser}
%   \caption{Seattle Mariners at Spring Training, 2010.}
%   \Description{Enjoying the baseball game from the third-base
%   seats. Ichiro Suzuki preparing to bat.}
%   \label{fig:teaser}
% \end{teaserfigure}

% \received{20 February 2007}
% \received[revised]{12 March 2009}
% \received[accepted]{5 June 2009}

%%
%% This command processes the author and affiliation and title
%% information and builds the first part of the formatted document.
\maketitle

\section{Introduction}
Interpretable machine learning is often motivated by settings in which model behavior must be inspected, challenged, or communicated to stakeholders with different technical backgrounds. Rule-based models are especially attractive in this setting because they expose explicit prediction-time evidence in the form of matched conditions and predicted outcomes. HEROS is a rule-based machine learning system that separates optimization of individual rules from optimization of whole rule sets, producing compact models together with instance-level rule evidence~\cite{lipschutzvilla2025heros}. More broadly, HEROS belongs to a longer learning-classifier-system and rule-based machine learning lineage that has emphasized interpretability, heterogeneous-pattern discovery, and biomedical applicability over many years~\cite{urbanowicz2009lcs,urbanowicz2017intro,urbanowicz2015exstracs}. Nonetheless, rule outputs that are intelligible to machine learning researchers are not automatically appropriate for clinicians, patients, or other non-expert users.

Large language models (LLMs) seem well suited to this communication problem, but recent work on natural-language explanations shows that fluent text can be only weakly coupled to the reasoning process that produced a prediction~\cite{atanasova2023faithfulness,madsen2024self,parcalabescu2024measuring,lyu2024survey}. In high-stakes settings, that gap is especially problematic because readers may mistake an explanation of model behavior for an explanation of the world itself~\cite{carriero2025describe}. These concerns motivate a narrower design goal than generic explanation generation. Instead of asking an LLM to infer reasons from raw data or to rationalize a black-box model after the fact, we ask it to translate evidence that already comes from an interpretable model.

This framing is related to a growing body of hybrid work that combines LLMs with symbolic or rule-based structures instead of leaving reasoning fully unconstrained. Recent examples include symbolic chain-of-thought formulations for faithful logical reasoning~\cite{xu2024symbcot}, rule-based decomposition for legal reasoning~\cite{servantez2024chainlogic}, symbolic working memory for multi-step rule application~\cite{wang2024symbolicwm}, and the use of LLMs to support formal constraint acquisition~\cite{mechqrane2024llmacq}. These studies show that explicit rules and symbolic scaffolds can improve controllability, traceability, and logical fidelity. Our setting differs in an important way: we do not use rules primarily to help the LLM solve a reasoning task. Instead, we use an interpretable rule-based classifier to produce the explanation evidence first, and then ask the LLM to translate that evidence for different audiences.

We use the term \emph{LLM} in the broad sense of large autoregressive language models capable of instruction following and long-form text generation, consistent with the current GPT family represented by systems such as GPT-4~\cite{openai2023gpt4}. In this study, however, the relevant design question is not whether the model can reason freely over raw task inputs, but whether it can reliably restate structured rule evidence under narrow behavioral constraints. That distinction informs the specific model choice described in the Methods section.

Our empirical question has three linked parts. First, if evidence extraction remains in the hands of an interpretable rule-based model and the language model is used only for audience-conditioned translation, then the dominant failure mode should shift away from unsupported invention and toward selective omission or over-compression of valid evidence. Second, explanation quality should depend not only on prompt design, but also on the quality of the learned rule model that supplies the packet in the first place. Third, audience customization should produce distinct trade-offs rather than a single uniformly best output style. We therefore vary the HEROS training profile, the prompt experiment, and the target audience while holding the downstream translation task fixed. For each held-out instance, the LLM receives only structured evidence derived from matching rules in the selected HEROS model, instance feature values, a model-context summary, and optional glossary entries. It is explicitly instructed to describe model behavior rather than domain truth and to avoid causal claims.

This paper reports a preliminary six-run comparison on MUX6, MUX11, and MUX20 rather than a definitive end-state evaluation. The central empirical claim is that explanation quality in this setting is jointly determined by three factors: the quality of the rule-based model that supplies the evidence packet, the prompt design that governs how that packet may be verbalized, and the audience target for which the verbalization is produced. Across all six runs, \emph{Evidence Precision} and \emph{Prediction--Explanation Agreement} remain perfect, but the remaining metrics vary meaningfully across these three axes. MUX6 is ideal under both training profiles, MUX11 improves sharply under extensive training, MUX20 remains difficult, the glossary-style prompt generally improves audience-fit scores while lengthening text, and the audience settings separate into high-recall, low-hallucination, and high-overclaim regimes. Because this is a preliminary paper, we do not treat these patterns as settled properties of rule-based explanation pipelines. Instead, we use them to characterize the interaction among upstream model quality, prompt design, and audience targeting in a first end-to-end study of constrained LLM translation for rule-based evidence.

\section{Methods}
The study design follows directly from the three-factor framing introduced above. HEROS supplies the explanation evidence so that the generator never has to infer explanation content from raw inputs alone, prompting restricts the LLM to audience-conditioned translation rather than open-ended reasoning, and evaluation separates unsupported invention from incomplete evidence preservation. The remainder of this section presents those components in the same order: first the rule-based model and its prediction mechanism, then the LLM configuration and prompt design, and finally the benchmark and evaluation setup used to test the resulting pipeline.

\subsection{HEROS and rule-based machine learning context}
Interpretable rule-based machine learning spans several related traditions. Classical rule-induction systems such as RIPPER learn compact ordered IF--THEN rules directly from labeled data~\cite{cohen1995ripper}, while evolutionary rule learners such as BioHEL search for readable rule sets with genetic optimization~\cite{bacardit2006biohel}. Within learning classifier systems, XCS and UCS helped establish the accuracy-based view of rule learning in which predictions are made from matched rules and fitness emphasizes accurate, general classifiers rather than only raw reward or coverage~\cite{wilson2000xcs,bernado2003ucs}. More recent supervised Michigan-style systems such as ExSTraCS adapted this lineage to noisy, heterogeneous biomedical settings while preserving an explicit rule representation~\cite{urbanowicz2015exstracs,urbanowicz2017intro}. HEROS should be understood against that broader rule-based background rather than as an isolated system.

HEROS is most directly distinguished by its two-phase training design~\cite{lipschutzvilla2025heros}. In Phase 1, HEROS evolves an interpretable rule population and evaluates individual rules with a Pareto-inspired objective that favors accurate rules that cover informative subsets of the training instances. In Phase 2, HEROS treats candidate predictive models as subsets of those learned rules and evolves rule sets under a separate multi-objective search that prioritizes predictive performance while penalizing unnecessarily large models. The resulting output is therefore not just a flat population of rules, but a population of candidate rule sets together with the underlying rules from which they are assembled. At prediction time, one selected rule set is applied to a new instance, the rules in that model whose conditions match the instance form the match set, and those matching rules generate the prediction array used for class prediction. This separation between rule learning and rule-set learning is what makes HEROS useful for the present study: it exposes a compact predictive model at the model level and instance-specific matched-rule evidence at the prediction level. Table~\ref{tab:heros-settings} summarizes the shared HEROS settings used in the MUX6, MUX11, and MUX20 experiments, taken from the resolved experiment configuration and the corresponding HEROS implementation. We ran two HEROS training profiles. The \emph{minimal} profile reflects the lighter budget used in the original script-first prototype. The \emph{extensive} profile increases the HEROS search budget to the setting that recovers the ideal eight-rule MUX6 model in our reruns, and we then apply that same stronger profile unchanged to MUX11 and MUX20. In all runs, \texttt{Group} and \texttt{InstanceID} are excluded from model training and retained only for split management and logging.

\begin{table}[tbp]
\caption{Major HEROS training and model-selection settings used across MUX6, MUX11, and MUX20.}
\label{tab:heros-settings}
\centering
\small
\setlength{\tabcolsep}{3pt}
\renewcommand{\arraystretch}{1.05}
\begin{tabular}{p{0.37\columnwidth}p{0.22\columnwidth}p{0.24\columnwidth}}
\toprule
\textbf{Setting} & \textbf{Minimal} & \textbf{Extensive} \\
\midrule
Outcome type & \texttt{class} & \texttt{class} \\
Excluded columns & \texttt{Group}, \texttt{InstanceID} & \texttt{Group}, \texttt{InstanceID} \\
Rule-learning iterations & 10{,}000 & 100{,}000 \\
Rule population size & 500 & 1{,}000 \\
Model-level iterations & 40 & 100 \\
Model population size & 100 & 100 \\
Crossover probability & 0.8 & 0.8 \\
Mutation probability & 0.04 & 0.04 \\
Fitness function & \texttt{pareto} & \texttt{pareto} \\
Subsumption & \texttt{both} & \texttt{both} \\
Compaction & \texttt{sub} & \texttt{sub} \\
Random seed & 42 & 42 \\
Target model selection & \texttt{auto\_select\_top\_model} & \texttt{auto\_select\_top\_model} \\
\bottomrule
\end{tabular}
\end{table}

\subsection{From a HEROS model to a prediction}
The explanation pipeline begins only after HEROS has been trained. For each benchmark, HEROS is fit on the benchmark's training split, producing a learned rule population and a learned population of candidate rule sets. The explanation system then selects one trained rule set to act as the predictive model and uses only that model when explaining held-out instances. This ordering matters because the LLM is not asked to explain an idealized multiplexer program or the entire evolved rule population. It is asked to translate the evidence produced by one selected learned HEROS model.

In HEROS, a predictive model is itself a rule set rather than a single parametric object. Phase 1 evolves individual rules, and Phase 2 evolves candidate models by selecting subsets of those rules and evaluating them as whole rule sets~\cite{lipschutzvilla2025heros}. At prediction time, HEROS selects one model and treats that model's rule set as the prediction mechanism. In implementation terms, the selected model exposes a \texttt{target\_rule\_set}, and prediction begins by scanning that rule set to construct the match set $M$, namely the subset of rules whose IF conditions match the current instance. HEROS then aggregates those matching rules into a class-wise prediction array.

For classification, each matching rule $r \in M$ contributes its empirical class distribution $p_r(c)$ together with its numerosity $n_r$. HEROS computes the unnormalized class vote
\[
V(c) = \sum_{r \in M} n_r p_r(c),
\]
then normalizes by the total numerosity in the match set,
\[
P(c) = \frac{V(c)}{\sum_{r \in M} n_r}.
\]
The predicted class is the class with the largest $P(c)$. If there is a tie, HEROS breaks it using the training majority among the tied classes, and if that is also tied it falls back to random selection. If no rules match, the instance is uncovered and HEROS uses its default uncovered-case behavior. Table~\ref{tab:mux6-ideal-model} gives a concrete MUX6 illustration of what such a model looks like and how the match set and prediction array are formed.

\begin{table*}[t]
\caption{Conceptual MUX6 example showing what a rule-based model looks like and how prediction is formed from a match set. For this idealized MUX6 model, the two address bits select among $R_0$, $R_1$, $R_2$, and $R_3$ in the order 00, 01, 10, 11. Each ideal rule has numerosity 1 and a deterministic class distribution, so the example prediction array is especially simple. Learned HEROS models are usually more general and may yield larger match sets and softer class probabilities.}
\label{tab:mux6-ideal-model}
\centering
\small
\setlength{\tabcolsep}{4pt}
\renewcommand{\arraystretch}{1.06}
\begin{tabular}{cll}
\toprule
\textbf{Rule} & \textbf{IF} & \textbf{THEN} \\
\midrule
$r_1$ & $A_0=0 \wedge A_1=0 \wedge R_0=0$ & predict class 0 \\
$r_2$ & $A_0=0 \wedge A_1=0 \wedge R_0=1$ & predict class 1 \\
$r_3$ & $A_0=0 \wedge A_1=1 \wedge R_1=0$ & predict class 0 \\
$r_4$ & $A_0=0 \wedge A_1=1 \wedge R_1=1$ & predict class 1 \\
$r_5$ & $A_0=1 \wedge A_1=0 \wedge R_2=0$ & predict class 0 \\
$r_6$ & $A_0=1 \wedge A_1=0 \wedge R_2=1$ & predict class 1 \\
$r_7$ & $A_0=1 \wedge A_1=1 \wedge R_3=0$ & predict class 0 \\
$r_8$ & $A_0=1 \wedge A_1=1 \wedge R_3=1$ & predict class 1 \\
\bottomrule
\end{tabular}

\vspace{0.6em}

\begin{tabular}{p{3.4cm}p{7.1cm}p{3.7cm}}
\toprule
\textbf{Example instance} & \textbf{Match set and prediction array} & \textbf{Output} \\
\midrule
$A_0=1$, $A_1=0$, $R_0=1$, $R_1=0$, $R_2=1$, $R_3=1$ &
Only $r_6$ matches, so $M=\{r_6\}$. With numerosity $n_{r_6}=1$ and deterministic class distribution $p_{r_6}(0)=0$, $p_{r_6}(1)=1$, the vote totals are $V(0)=1\times 0=0$ and $V(1)=1\times 1=1$. The total numerosity is $\sum_{r \in M} n_r = 1$, so the normalized
prediction array is $P(0)=0/1=0.0$ and $P(1)=1/1=1.0$. &
Covered by the model: yes. Predicted class: 1. Prediction probabilities: $\{0: 0.0,\ 1: 1.0\}$. \\
\bottomrule
\end{tabular}
\end{table*}

\subsection{LLM and ChatGPT model selection}
All reported generations in this study were produced through the OpenAI API with \texttt{gpt-4.1-mini}, rather than through the ChatGPT web interface. ChatGPT is referenced here only to distinguish interactive product use from programmatic API-based experimentation~\cite{openai2025chatgptmodelselector}. The experiments themselves were conducted through fully logged API calls so that prompts, parameters, and outputs remained auditable and rerunnable.

The choice of \texttt{gpt-4.1-mini} was deliberate. OpenAI describes the GPT-4.1 family as improving instruction following and long-context performance relative to earlier GPT-4-class API models, with GPT-4.1 mini positioned as a lower-latency, lower-cost member of that family~\cite{openai2025gpt41,openai2025gpt41mini}. Those characteristics fit the present task well. Our prompts are structured, bounded, and evidence-conditioned; they ask the model to translate supplied rules and feature values rather than solve open-ended reasoning problems. The study design also requires many repeated calls across datasets, experimental setups, audiences, and judge passes. Under that workload, a smaller instruction-following model is preferable to a frontier-sized model because it keeps latency and cost manageable while still supporting deterministic generation with sufficiently strong compliance to the prompt constraints. We therefore selected \texttt{gpt-4.1-mini} as a practical balance between controllability, throughput, and experimental scale, while recognizing that future work should compare stronger and weaker model families directly.

\begin{table}[tbp]
\caption{LLM explanation and evaluation run parameters used in the final study.}
\label{tab:run-params}
\centering
\small
\setlength{\tabcolsep}{3pt}
\renewcommand{\arraystretch}{1.05}
\begin{tabular}{p{0.43\columnwidth}p{0.47\columnwidth}}
\toprule
\textbf{Setting} & \textbf{Value} \\
\midrule
Prompt version & \texttt{v2\_minimal} \\
Training profiles & Minimal; Extensive \\
Experiments & 1 and 2 \\
Audiences & Layman, Clinician, Expert \\
Sampling strategy & Stratified by prediction and rule-match bucket, with the full MUX6 test fold used when available \\
Generation model & \texttt{gpt-4.1-mini} \\
Generation temperature & 0.0 \\
Generation max tokens & 400 \\
Judge model & \texttt{gpt-4.1-mini} \\
Judge scoring & Applied to all completed runs; post hoc for minimal-profile runs \\
Judge temperature & 0.0 \\
Judge max tokens & 300 \\
Prompt metadata exposure & Clinician and Expert only \\
Glossary exposure & Experiment 2 only; minimal feature-only glossary \\
\bottomrule
\end{tabular}
\end{table}

\subsection{Pipeline design}
The pipeline is built around a strict division of labor between HEROS and the language model. HEROS performs prediction and identifies the matching rules for a given instance. The LLM receives only a structured explanation packet derived from those outputs. It is never asked to infer feature importance from scratch, summarize the full rule population, or recover dataset semantics from internal feature indices.

For each held-out instance, the system assembles a packet containing the predicted class, class probabilities, model coverage status, matching rules partitioned into supporting and contradictory rules, relevant instance feature values, and a model-context summary describing the number of matching rules and whether the evidence is consistent or mixed. When available, the packet also includes rule metadata such as numerosity, accuracy, fitness, vote contribution, and training-support counts. Internal feature indices are converted to feature names before prompt construction. Each matching rule is then rendered as readable IF--THEN text so that the generator receives a textual but still structured view of the model evidence.

The prompt design follows the same constrained translation philosophy. We compare two experimental setups. \emph{Experiment 1} provides prediction output, matching rules, and relevant instance values. \emph{Experiment 2} adds concise one-sentence glossary entries for features that appear in the matching rules. These glossary entries were intentionally minimal and feature-oriented rather than mechanism-oriented. For synthetic multiplexer benchmarks, a detailed glossary that explains the true MUX mechanism could leak benchmark structure into the prompt. We therefore restricted the glossary to simple feature descriptions rather than definitions of how the multiplexer operates. The audience setting determines how much technical detail is exposed. Layman explanations use plain language with qualitative confidence wording only. Clinician explanations use concise feature--value reasoning with moderate technical detail and rounded support information. Expert explanations include exact probability language, rule identifiers, and metadata when available. Across all settings, the prompt instructs the LLM to discuss model behavior only, to acknowledge conflicting rules when they are present, and to avoid causal or factual claims that go beyond the provided evidence.

\subsection{Benchmarks and experimental design}
We evaluate on the multiplexer benchmarks distributed with the HEROS repository. MUX6 contains 500 total instances split into 450 training and 50 held-out test cases. MUX11 contains 5{,}000 total instances split into 4{,}500 training and 500 held-out test cases. MUX20 contains 10{,}000 total instances split into 9{,}000 training and 1{,}000 held-out test cases. HEROS is trained on the standard training split for each benchmark. For MUX6, we evaluate on the full 50-instance held-out test fold. For MUX11 and MUX20, we evaluate on stratified 50-instance subsets drawn from the held-out test folds in order to keep the LLM evaluation budget fixed across datasets while preserving variation in predicted class and rule-matching profile.

We run both HEROS training profiles on each benchmark. This design yields 50 evaluated instances per dataset, two training profiles, two experimental setups, and three audience targets, for a total of 1{,}800 explanations. All reported results therefore compare \emph{Experiment 1} and \emph{Experiment 2} under both the \emph{minimal} and \emph{extensive} HEROS training profiles.
Table~\ref{tab:benchmarks} summarizes the benchmark sizes, while Tables~\ref{tab:heros-settings} and~\ref{tab:run-params} summarize the major HEROS settings and the explanation run parameters used throughout. We evaluate the LLM against the evidence packet produced by the learned HEROS model rather than against the ideal benchmark logic. That choice is intentional because the downstream explanation task is to restate what the trained model actually did at prediction time. The comparison between HEROS profiles lets us measure how much that upstream model quality changes the downstream translation problem.

The generation model was \texttt{gpt-4.1-mini} with temperature fixed at 0.0. Judge scoring used the same model family at temperature 0.0 under a separate JSON-only rubric. We used deterministic settings throughout to reduce sampling variance and to support rerunnable, fully logged API-based experiments. We did not use the ChatGPT interface itself for experimental generation; all model calls were issued through the API to keep prompts, parameters, and outputs explicitly logged. For the minimal HEROS training profile, judge scoring was applied post hoc so that the generation configuration remained otherwise identical to the original prototype profile.

\begin{table}[H]
\caption{Benchmark configuration. Each benchmark contributes 50 evaluated held-out instances under two HEROS training profiles, producing 600 generated explanations per dataset and 1{,}800 explanations overall.}
\label{tab:benchmarks}
\centering
\small
\setlength{\tabcolsep}{4pt}
\begin{tabular}{lrrr}
\toprule
\textbf{Dataset} & \textbf{Train} & \shortstack{\textbf{Held-out}\\\textbf{Test}} & \shortstack{\textbf{Evaluated per}\\\textbf{Training Profile}} \\
\midrule
MUX6  & 450  & 50   & 50 \\
MUX11 & 4{,}500 & 500  & 50 \\
MUX20 & 9{,}000 & 1{,}000 & 50 \\
\bottomrule
\end{tabular}
\end{table}

\subsection{Evaluation Metrics}
The evaluation framework is designed to test the study design directly. We separate the metrics into three groups. The first group contains packet-computable metrics tied to prior evidence-attribution, rationale-evaluation, hallucination, and readability literature. The second group contains task-specific diagnostics introduced in this paper because our setting has two unusual constraints: the explanation must stay tied to a selected rule-based model's matching rules, and it must describe model behavior without drifting into causal or real-world claims. The third group contains audience-facing judge scores produced by an LLM. Table~\ref{tab:metrics} gives the exact operationalization, value range, preferred direction, and literature basis or justification for each metric. We intentionally place those details in the table so the prose can remain concise.

These judge scores are included only as weak, scalable proxies for audience fit and should be interpreted cautiously rather than as human-grounded evidence in the sense of Doshi-Velez and Kim~\cite{doshivelez2017rigorous}. In particular, we do recognize that using an LLM to evaluate LLM-generated text is a methodological limitation, especially because the generator and judge come from the same model family~\cite{liu2023geval,huang2025judge}. For that reason, we treat the packet-computable metrics as primary and the judge-model outputs as secondary. We also do not report comprehensiveness or sufficiency, although both are well established in rationale evaluation~\cite{deyoung2020eraser}, because their proper computation here requires rerunning HEROS under feature-removal and keep-only perturbations rather than estimating them from the text alone.

\begin{table*}[t]
\caption{Metrics used in the study. Rows explicitly marked as task-specific are operationalizations introduced in this paper. Judge metrics are rubric-based LLM outputs and should be interpreted as weak supportive proxies rather than direct human-study outcomes.}
\label{tab:metrics}
\centering
\small
\setlength{\tabcolsep}{3pt}
\renewcommand{\arraystretch}{1.06}
\begin{tabular}{p{2.0cm}p{2.8cm}p{5cm}p{2.5cm}p{3.1cm}}
\toprule
\textbf{Category} & \textbf{Metric} & \textbf{Calculation or operationalization} & \textbf{Range / direction} & \textbf{Basis} \\
\midrule
Evidence grounding & Evidence Precision & $\lvert M \cap E \rvert / \lvert M \rvert$, where $M$ is the set of mentioned evidence features and $E$ is the set of packet evidence features derived from the instance's matching rules. & $[0,1]$; higher is better. & Evidence attribution and rationale evaluation~\cite{deyoung2020eraser,xing2025evidence}. \\
Evidence grounding & Evidence Recall & $\lvert M \cap E \rvert / \lvert E \rvert$. Captures how much packet evidence is preserved in the explanation. & $[0,1]$; higher is better. & Evidence attribution and rationale evaluation~\cite{deyoung2020eraser,xing2025evidence}. \\
Evidence grounding & Evidence F1 & Harmonic mean of \emph{Evidence Precision} and \emph{Evidence Recall}. & $[0,1]$; higher is better. & Standard summary of precision and recall~\cite{deyoung2020eraser}. \\
Faithfulness & \emph{Prediction--Explanation Agreement} & Mean of $A_i$, where $A_i=1$ if the explanation either mentions the predicted class or mentions no class label explicitly, and $A_i=0$ otherwise. & $[0,1]$; higher is better. & Task-specific metric introduced in this paper. \\
Faithfulness & \emph{Rule Coverage} & For top-$k$ supporting matching rules $K_i$, computes $\frac{1}{|K_i|}\sum_{r \in K_i}\mathbf{1}[F(r)\cap M_i \neq \emptyset]$. Feature-overlap proxy only; not exact rule reconstruction. & $[0,1]$; higher is better. & Task-specific coverage proxy introduced in this paper. \\
Hallucination & Hallucination Rate & Mean of $H_i$, where $H_i=1$ if an explanation contains unsupported features or unsupported reality claims. & $[0,1]$; lower is better. & Hallucination construct~\cite{bang2025hallulens}; implementation adapted to this pipeline. \\
Safety & \emph{Uncertainty Acknowledgment Rate} & On packets labeled mixed, mostly supporting, or uncovered, mean indicator for predefined uncertainty cues such as ``mixed,'' ``uncertain,'' or ``may.'' & $[0,1]$; higher is better when uncertainty is warranted. & Task-specific metric introduced in this paper. \\
Safety & \emph{Conflict Acknowledgment Score} & Binary consistency check that equals 1 if and only if conflict language appears exactly when the packet contains conflicting matching rules. & $[0,1]$; higher is better. & Task-specific metric introduced in this paper. \\
Safety & \emph{Causal Overclaim Rate} & Mean of a binary flag for causal phrases or reality-claim phrases that go beyond packet evidence. & $[0,1]$; lower is better. & Task-specific safety metric introduced in this paper. \\
Audience quality & \emph{Audience Understandability} & Rubric-based LLM-judge score for whether the explanation is understandable for the requested audience. & $[0,1]$; higher is better, but interpret cautiously. & Judge-based proxy motivated by XAI and LLM evaluation work~\cite{kim2024human,liu2023geval,huang2025judge}. \\
Audience quality & \emph{Audience Technical Fit} & Rubric-based LLM-judge score for whether the technical detail matches the requested audience. & $[0,1]$; higher is better, but interpret cautiously. & Judge-based proxy motivated by XAI and LLM evaluation work~\cite{kim2024human,liu2023geval,huang2025judge}. \\
Readability & Flesch Reading Ease & Standard readability formula computed for all audiences. & Higher means easier text; roughly 90--100 very easy, 60--70 standard, 30--50 difficult. & Readability formula~\cite{flesch1948readability}. \\
Readability & Flesch--Kincaid Grade Level & Standard readability formula computed for all audiences. & Lower means simpler text; roughly 6 middle-school, 9--10 high-school, above 12 college-level. & Readability formula~\cite{kincaid1975derivation}. \\
Output length & Word Count & Tokenized word count per explanation. & $\geq 0$; descriptive only. & Descriptive analysis choice in this paper. \\
\bottomrule
\end{tabular}
\end{table*}

\section{Results}
\subsection{HEROS model quality and cross-benchmark performance}
Table~\ref{tab:dataset-results} reports the six-run comparison across datasets and HEROS training profiles. Across all six runs, \emph{Evidence Precision} and \emph{Prediction--Explanation Agreement} remained perfect and are therefore omitted from the table to save space. The main purpose of this section is to trace the three-factor argument introduced in the Introduction through the empirical results. The dataset/profile comparison establishes how HEROS training profile changes the evidence packet that must be translated. Table~\ref{tab:pooled-trends} then shows how prompt experiment and audience target reshape the resulting outputs once that packet exists.

Taken together, the results support three claims. First, HEROS training profile changes the quality and density of the evidence packet, and that change materially affects downstream explanation quality. Second, prompt experiment changes how much the LLM elaborates on the packet, typically trading modest audience-fit gains for longer text and more room for unsupported additions. Third, audience target changes which failure mode becomes most visible: expert outputs preserve the most evidence, layman outputs are the safest under the hallucination and overclaim checks, and clinician outputs remain the most causally overstated.

Averaged across datasets, the \emph{extensive} profile raises mean HEROS test accuracy from 80.4\% to 87.3\%, increases the ideal-rule fraction in the selected model from 39.6\% to 62.5\%, and lowers mean matching rules from 6.07 to 5.37. Those averages are informative, but Table~\ref{tab:dataset-results} also shows that they hide strong dataset-specific differences in how upstream model quality changes the downstream translation task.

MUX6 provides the simplest case. Both HEROS profiles recover an ideal 8-rule model, achieve HEROS test accuracy 1.000, and average exactly one matching rule per held-out instance. Under those conditions, downstream explanation coverage is perfect in both runs: \emph{Evidence Recall} and \emph{Rule Coverage} remain 1.000. The remaining differences are modest wording effects rather than evidence differences. The \emph{minimal} profile yields a slightly lower \emph{Hallucination Rate}, 16.0\% versus 18.3\%, and slightly higher audience scores, but the two MUX6 runs are substantively similar because the learned rule model itself is already ideal. MUX6 therefore acts as the clean-control case for the rest of the paper: once the packet is already minimal, the remaining variation mostly reflects language-level choices. Consistent with that interpretation, \emph{Causal Overclaim Rate} remains 0.0 in both MUX6 runs.

MUX11 shows the strongest benefit of stronger HEROS training. Moving from the \emph{minimal} profile to the \emph{extensive} profile raises HEROS test accuracy from 81.2\% to 98.8\%, reduces the selected top model from 56 rules to 24 rules, increases the ideal-rule fraction from 18.8\% to 87.5\%, and lowers the mean number of matching rules from 4.86 to 2.24. These upstream changes coincide with substantial downstream gains: \emph{Evidence Recall} rises from 94.4\% to 99.4\%, \emph{Rule Coverage} reaches 100.0\%, \emph{Causal Overclaim Rate} falls from 17.0\% to 9.0\%, \emph{Audience Understandability} and \emph{Audience Technical Fit} both improve slightly, and outputs become markedly shorter, from 126.11 words to 102.73. The one countervailing result is that \emph{Hallucination Rate} rises from 5.7\% to 18.3\%. This pattern suggests that better trained HEROS models can simplify the packet and improve evidence preservation while still giving the LLM more room to elaborate beyond the packet in its phrasing. Even with that increase in hallucination, MUX11 remains the clearest positive result in the study because training quality, evidence preservation, overclaim control, and output compactness all move in the same overall direction.

MUX20 provides the strongest counterexample. Stronger HEROS training raises HEROS test accuracy only slightly, from 60.1\% to 63.2\%, recovers no ideal rules, and does not simplify the packet. Mean matching rules increase slightly from 12.34 to 12.88, \emph{Evidence Recall} falls from 73.7\% to 68.6\%, and \emph{Rule Coverage} falls from 94.3\% to 93.0\%. At the same time, \emph{Hallucination Rate} drops from 12.0\% to 8.7\% and the audience-oriented scores improve slightly, but \emph{Causal Overclaim Rate} rises from 32.0\% to 34.0\%. The stronger HEROS profile therefore helps MUX20 only modestly at the classifier level and does not translate into better evidence preservation or safer wording overall. MUX20 is therefore the clearest boundary case in the study: upstream training can improve classifier accuracy without making the explanation packet easier to translate. Taken together, the six-run comparison shows that better HEROS training can materially improve downstream explanation quality, but only when it also produces a cleaner and more ideal rule model.

\begin{table*}[t]
\caption{Overall results by dataset and HEROS training profile. \emph{Evidence Precision} and \emph{Prediction--Explanation Agreement} were 100.0 for every row and are omitted to save space. Judge scores are on a 0--1 scale and should be interpreted as weak proxy measures rather than as human-study outcomes. Percentage-style metrics are shown as percentages. Abbreviations: HEROS Acc. = HEROS test accuracy; Top Rules = number of rules in the selected HEROS model; Ideal \% = percentage of rules in the selected model that match the ideal MUX rule set; Ev. Recall = \emph{Evidence Recall}; Halluc. = \emph{Hallucination Rate}; Rule Cov. = \emph{Rule Coverage}; COR = \emph{Causal Overclaim Rate}; Aud. Und. = \emph{Audience Understandability}; Aud. Tech. Fit = \emph{Audience Technical Fit}. \emph{Word Count} is averaged over all explanations within each row.}
\label{tab:dataset-results}
\centering
\small
\setlength{\tabcolsep}{4pt}
\renewcommand{\arraystretch}{1.08}
\resizebox{\textwidth}{!}{%
\begin{tabular}{llrrrrrrrrrrr}
\toprule
\textbf{Dataset} & \shortstack{\textbf{Training}\\\textbf{Profile}} & \textbf{HEROS Acc.} & \textbf{Top Rules} & \textbf{Ideal \%} & \textbf{Match Rules} & \textbf{Ev. Recall} & \textbf{Halluc.} & \textbf{Rule Cov.} & \textbf{COR} & \textbf{Aud. Und.} & \textbf{Aud. Tech. Fit} & \textbf{Words} \\
\midrule
MUX6  & Minimal   & 100.0 & 8  & 100.0 & 1.00  & 100.0 & 16.0 & 100.0 & 0.0  & 0.902 & 0.873 & 90.78 \\
MUX6  & Extensive & 100.0 & 8  & 100.0 & 1.00  & 100.0 & 18.3 & 100.0 & 0.0  & 0.888 & 0.868 & 90.42 \\
MUX11 & Minimal   & 81.2  & 56 & 18.8  & 4.86  & 94.4  & 5.7  & 99.9  & 17.0 & 0.827 & 0.833 & 126.11 \\
MUX11 & Extensive & 98.8  & 24 & 87.5  & 2.24  & 99.4  & 18.3 & 100.0 & 9.0  & 0.842 & 0.846 & 102.73 \\
MUX20 & Minimal   & 60.1  & 54 & 0.0   & 12.34 & 73.7  & 12.0 & 94.3  & 32.0 & 0.796 & 0.818 & 139.74 \\
MUX20 & Extensive & 63.2  & 69 & 0.0   & 12.88 & 68.6  & 8.7  & 93.0  & 34.0 & 0.813 & 0.829 & 141.01 \\
\bottomrule
\end{tabular}%
}
\end{table*}

\subsection{Illustrative MUX11 example}
Table~\ref{tab:mux11-example} provides a qualitative example that connects the three empirical axes within a single held-out case. We fix the prompt to \emph{Experiment 1} and use \texttt{InstanceID 287} from MUX11 because the same instance yields different HEROS packets under the two training profiles while still producing readable layman, clinician, and expert outputs within each profile. We use \emph{Experiment 1} rather than \emph{Experiment 2} because, for this instance, it yields the clearer audience-separated error pattern: the \emph{minimal / expert} row is flagged by the hallucination detector, whereas the other five rows are not. Under the \emph{minimal} profile, the selected HEROS model predicts \textbf{class 1} with probability \textbf{0.665} from three supporting matching rules and no conflict. Under the \emph{extensive} profile, the selected model predicts \textbf{class 0} with probability \textbf{0.775} from one strong supporting rule plus two contradictory matching rules.

The table is intended to be read in two directions. Reading across profiles shows how HEROS training can alter the packet itself for the same held-out case. Reading within either profile shows how audience targeting reshapes the same packet into different verbal forms. The layman rows compress the packet into short qualitative summaries, the clinician rows preserve concise training-support language and moderated confidence phrasing, and the expert rows expose exact probabilities, rule identifiers, and rule-level detail. In the \emph{minimal} block, only the expert explanation is flagged by the hallucination detector. That pattern illustrates a recurring trade-off observed elsewhere in the paper: richer technical detail can preserve more of the packet, but it can also create more room for unsupported elaboration.

\begin{table*}[t]
\caption{Illustrative MUX11 held-out example (\texttt{InstanceID 287}) under \emph{Experiment 1}. The same held-out instance yields different HEROS packets under the two training profiles. \emph{Minimal}: predicts \textbf{class 1} with probability \textbf{0.665} from three supporting matching rules and no conflict. \emph{Extensive}: predicts \textbf{class 0} with probability \textbf{0.775} from three matching rules, of which one supports the prediction and two contradict it. Column \textbf{Hall.} reports whether the paper's hallucination detector flagged the row. Boldface is added for emphasis.}
\label{tab:mux11-example}
\centering
\scriptsize
\renewcommand{\arraystretch}{1.08}
\setlength{\tabcolsep}{4pt}
\begin{tabular}{p{1.8cm}p{1.3cm}p{0.8cm}r p{9.9cm}}
\toprule
\textbf{Training Profile} & \textbf{Audience} & \textbf{Hall.} & \textbf{Words} & \textbf{Output} \\
\midrule
\multicolumn{5}{p{15.8cm}}{\textbf{Minimal packet summary:} The selected HEROS model predicts \textbf{class 1} with probability \textbf{0.665}. Three matching rules support the prediction and no conflicting matching rules are present. The strongest recurring feature pattern is \textbf{$A_1 = 0$} with \textbf{$R_1 = 1$}, together with \textbf{$A_2 = 0$} and either \textbf{$R_4 = 1$} or \textbf{$R_5 = 1$}.} \\
\midrule \emph{Minimal} & Layman & No & 86 & ``The model leaned toward \textbf{class 1} because several conditions in the data all support this prediction. Specifically, the values show that \textbf{A\_2 = 0}, and combinations of \textbf{A\_1 = 0} with \textbf{R\_1 = 1} and \textbf{R\_4 = 1}, as well as \textbf{A\_0 = 0} and \textbf{A\_1 = 0} with \textbf{R\_1 = 1} and \textbf{R\_5 = 1}, all point to class 1. Since all the active rules agree, the model is confident in this choice. Keep in mind, this explanation reflects how the model made its decision, not a statement of fact.'' \\
\emph{Minimal} & Clinician & No & 127 & ``The prediction of \textbf{class 1} with moderate confidence (about \textbf{66\%}) is supported by three active rules that all agree on this outcome. One rule indicates that when \textbf{A\_2 = 0}, class 1 is predicted, though this rule alone has mixed support with about 50\% of similar cases in training being class 1. Two additional rules, which consider combinations of \textbf{A\_1}, \textbf{R\_1}, \textbf{R\_4}, and \textbf{R\_5} values, provide stronger support with around 74--75\% of matched training instances belonging to class 1. Specifically, the presence of \textbf{R\_1 = 1}, \textbf{R\_4 = 1}, and \textbf{R\_5 = 1} alongside \textbf{A\_0 = 0} and \textbf{A\_1 = 0} aligns with these more confident rules. Overall, the evidence pattern from these rules consistently supports the prediction, but this explanation reflects model reasoning rather than a definitive causal conclusion.'' \\
\emph{Minimal} & Expert & Yes & 164 & ``The prediction of \textbf{class 1} with moderate confidence (\textbf{probability 0.665}) is supported by three active rules, all of which favor class 1 without conflict. \textbf{Rule 8051} matches on \textbf{A\_2 = 0} and has a near-even class distribution in training (about 50\% class 1), contributing modestly to the vote. \textbf{Rules 7904} and \textbf{7831} provide stronger support: both require \textbf{A\_1 = 0} and \textbf{R\_1 = 1}, with Rule 7904 also requiring \textbf{R\_4 = 1} and Rule 7831 requiring \textbf{A\_0 = 0} and \textbf{R\_5 = 1}. These latter two rules matched fewer instances but had higher accuracy (about 74--75\%) and fitness, indicating more reliable support for class 1. The instance's feature values (\textbf{A\_0 = 0}, \textbf{A\_1 = 0}, \textbf{A\_2 = 0}, \textbf{R\_1 = 1}, \textbf{R\_4 = 1}, \textbf{R\_5 = 1}) satisfy all three rules' conditions, leading to unanimous agreement among the active rules. This consensus underpins the moderate confidence in the class 1 prediction. This explanation reflects the model's reasoning based on rule matches and training support, not a causal inference or ground truth statement.'' \\
\midrule
\multicolumn{5}{p{15.8cm}}{\textbf{Extensive packet summary:} The selected HEROS model predicts \textbf{class 0} with probability \textbf{0.775}. Three matching rules are present, of which one strongly supports the prediction and two contradict it. The decisive supporting rule uses \textbf{$A_0 = 0$}, \textbf{$A_1 = 0$}, \textbf{$A_2 = 0$}, and \textbf{$R_0 = 0$}; the conflicting rules additionally involve \textbf{$R_2 = 1$} together with either \textbf{$R_3 = 0$} or \textbf{$R_6 = 0$}.} \\
\midrule \emph{Extensive} & Layman & No & 87 & ``The model leaned toward \textbf{class 0} mainly because one rule matched that supports this prediction, based on the values of \textbf{A\_0 = 0}, \textbf{A\_1 = 0}, \textbf{A\_2 = 0}, and \textbf{R\_0 = 0}. However, two other rules also matched that suggest a different class, since they rely on \textbf{A\_0 = 0}, \textbf{A\_2 = 0}, \textbf{R\_2 = 1}, and either \textbf{R\_3 = 0} or \textbf{R\_6 = 0}. This means the evidence is mixed, but overall the model favored class 0. Keep in mind this explanation reflects how the model made its decision, not a definitive fact.'' \\
\emph{Extensive} & Clinician & No & 118 & ``The prediction of \textbf{class 0} is mainly supported by one strong rule that matches the instance's feature values \textbf{A\_0 = 0}, \textbf{A\_1 = 0}, \textbf{A\_2 = 0}, and \textbf{R\_0 = 0}, which had nearly 100\% class 0 support in training. However, two other active rules, which also match the instance's features including \textbf{R\_2 = 1} and \textbf{R\_3 = 0} or \textbf{R\_6 = 0}, predict class 1 with moderate support around 79\%. This creates some conflicting evidence, but since the strongest rule supports class 0 and most active rules lean that way, the model predicts class 0 with moderate confidence of about \textbf{77\%}. This explanation reflects the model's reasoning based on the matched rules and feature values, not a definitive statement of truth.'' \\
\emph{Extensive} & Expert & No & 142 & ``The prediction of \textbf{class 0} with moderate confidence (\textbf{probability 0.775}) is primarily supported by \textbf{Rule 88893}, which matches the instance features \textbf{A\_0 = 0}, \textbf{A\_1 = 0}, \textbf{A\_2 = 0}, and \textbf{R\_0 = 0}. This rule has strong support, perfect accuracy on 266 training instances, and contributes exclusively to class 0 votes. However, two conflicting rules (\textbf{90750} and \textbf{91623}) also match the instance, both predicting class 1 with moderate support and about 79\% accuracy on their respective training subsets. These conflicting rules match on \textbf{A\_0 = 0}, \textbf{A\_2 = 0}, \textbf{R\_2 = 1}, combined with \textbf{R\_6 = 0} or \textbf{R\_3 = 0}, and contribute predominantly to class 1 votes. The overall evidence is mixed, with the majority of active rules favoring class 0 but notable opposition from class 1 rules. This explanation reflects the model's reasoning based on rule matches and does not imply causal relationships or ground truth.'' \\
\bottomrule
\end{tabular}
\end{table*}

\subsection{Pooled experiment and audience trends under both HEROS training profiles}
Table~\ref{tab:pooled-trends} complements the dataset-level comparison in Table~\ref{tab:dataset-results} by pooling the two experiment setups and the three audience settings within each HEROS training profile. This view is easier to scan than the full dataset-by-dataset split, while still showing how \emph{Evidence Recall}, \emph{Hallucination Rate}, \emph{Causal Overclaim Rate}, the judge-based audience metrics, and the readability diagnostics shift with prompt setup and audience target.

At the experiment level, the glossary-style \emph{Experiment 2} defines the paper's prompt-design result axis. It behaves as a consistent trade-off rather than a uniform improvement. Under the \emph{minimal} profile, \emph{Experiment 2} raises \emph{Evidence Recall} from 88.6\% to 90.2\% and \emph{Rule Coverage} from 97.7\% to 98.4\%, but it also raises \emph{Hallucination Rate} from 10.2\% to 12.2\%, lengthens the outputs from 116.28 to 121.47 words, lowers \emph{Flesch Reading Ease}, and raises \emph{Flesch--Kincaid Grade Level}. The same pattern largely persists under the \emph{extensive} profile. There, \emph{Experiment 2} slightly improves \emph{Evidence Recall}, lowers \emph{Causal Overclaim Rate} from 15.8\% to 12.9\%, and nudges both \emph{Audience Understandability} and \emph{Audience Technical Fit} upward, but it again increases \emph{Hallucination Rate}, lowers \emph{Flesch Reading Ease}, raises \emph{Flesch--Kincaid Grade Level}, and produces longer explanations. Read together, these rows show that \emph{Experiment 2} usually produces more elaborate outputs and modest audience-fit gains, but only selective improvements in faithfulness.

The audience rows define the third result axis of the paper: distinct error profiles rather than a simple quality ordering. Across both HEROS training profiles, \emph{Layman} prompting produces the lowest \emph{Hallucination Rate}, the lowest \emph{Causal Overclaim Rate}, and the shortest outputs, but also the lowest \emph{Evidence Recall} and the weakest audience-fit scores. \emph{Clinician} prompting sits between those extremes on \emph{Evidence Recall} and audience fit, but its \emph{Causal Overclaim Rate} is the worst of the three audiences in both training profiles, reaching 32.3\% under the \emph{minimal} profile and 29.3\% under the \emph{extensive} profile. \emph{Expert} prompting yields the highest \emph{Evidence Recall} and the strongest \emph{Audience Understandability} and \emph{Audience Technical Fit} scores, but it also has by far the highest \emph{Hallucination Rate} and the longest outputs. The readability metrics sharpen that interpretation rather than replacing it: clinician outputs are the hardest by \emph{Flesch Reading Ease} and \emph{Flesch--Kincaid Grade Level}, while expert outputs score as easier by those formulas despite being more technical, confirming that the readability measures are best read as surface-form descriptors rather than as direct audience-fit validators.

The training profiles also shape these pooled trends. Moving from the \emph{minimal} to the \emph{extensive} profile slightly lowers the pooled \emph{Causal Overclaim Rate}, raises both audience-fit scores, shortens the outputs overall, and produces marginally easier text by the readability formulas. At the same time, pooled \emph{Evidence Recall} remains almost unchanged and pooled \emph{Hallucination Rate} rises. This mirrors the dataset-level result: better HEROS training helps most when it simplifies the underlying rule model, but its downstream effects are not captured by any single metric alone.

In that sense, the pooled table does not replace the dataset-level analysis; it clarifies it. The training-profile comparison identifies when the evidence packet becomes easier to translate, the experiment comparison shows how the glossary-style prompt changes elaboration behavior, and the audience comparison shows which failure mode each audience setting tends to expose.

\begin{table*}[t]
\caption{Pooled experiment and audience trends across datasets under both HEROS training profiles. Each experiment row aggregates 450 explanations within a training profile, and each audience row aggregates 300 explanations within a training profile. Judge scores are on a 0--1 scale and should be interpreted as weak proxy measures rather than as human-study outcomes. Percentage-style metrics are shown as percentages. Abbreviations: Ev. Recall = \emph{Evidence Recall}; Halluc. = \emph{Hallucination Rate}; Rule Cov. = \emph{Rule Coverage}; COR = \emph{Causal Overclaim Rate}; Aud. Und. = \emph{Audience Understandability}; Aud. Tech. Fit = \emph{Audience Technical Fit}; FRE = \emph{Flesch Reading Ease}; FKGL = \emph{Flesch--Kincaid Grade Level}.}
\label{tab:pooled-trends}
\centering
\small
\setlength{\tabcolsep}{4pt}
\renewcommand{\arraystretch}{1.08}
\resizebox{\textwidth}{!}{%
\begin{tabular}{lllrrrrrrrrr}
\toprule
\shortstack{\textbf{Training}\\\textbf{Profile}} & \textbf{Group} & \textbf{Setting} & \textbf{Ev. Recall} & \textbf{Halluc.} & \textbf{Rule Cov.} & \textbf{COR} & \textbf{Aud. Und.} & \textbf{Aud. Tech. Fit} & \textbf{FRE} & \textbf{FKGL} & \textbf{Words} \\
\midrule
Minimal   & Experiment & 1         & 88.6 & 10.2 & 97.7 & 16.4 & 0.841 & 0.841 & 56.75 & 10.67 & 116.28 \\
Minimal   & Experiment & 2         & 90.2 & 12.2 & 98.4 & 16.2 & 0.843 & 0.842 & 55.70 & 10.83 & 121.47 \\
Extensive & Experiment & 1         & 89.2 & 14.2 & 97.7 & 15.8 & 0.846 & 0.848 & 56.96 & 10.46 & 109.10 \\
Extensive & Experiment & 2         & 89.4 & 16.0 & 97.6 & 12.9 & 0.849 & 0.848 & 55.97 & 10.73 & 113.67 \\
\midrule
Minimal   & Audience   & Layman    & 86.6 & 2.7  & 97.1 & 4.3  & 0.780 & 0.691 & 54.78 & 10.82 & 82.97 \\
Minimal   & Audience   & Clinician & 88.5 & 3.0  & 97.9 & 32.3 & 0.839 & 0.876 & 52.60 & 11.91 & 121.51 \\
Minimal   & Audience   & Expert    & 93.0 & 28.0 & 99.2 & 12.3 & 0.906 & 0.958 & 61.30 & 9.53  & 152.16 \\
Extensive & Audience   & Layman    & 86.2 & 6.7  & 97.1 & 3.0  & 0.786 & 0.691 & 55.72 & 10.71 & 79.33 \\
Extensive & Audience   & Clinician & 88.4 & 3.0  & 97.2 & 29.3 & 0.855 & 0.881 & 52.47 & 11.77 & 114.51 \\
Extensive & Audience   & Expert    & 93.3 & 35.7 & 98.7 & 10.7 & 0.902 & 0.972 & 61.20 & 9.31  & 140.32 \\
\bottomrule
\end{tabular}%
}
\end{table*}

\section{Discussion}
The six-run model-quality comparison, together with the pooled experiment and audience analyses, sharpens the paper's central claim. A large language model can function as a constrained translator of rule-based evidence when that evidence is made explicit and the prompts are narrow. The strongest support for that claim is that \emph{Evidence Precision} and \emph{Prediction--Explanation Agreement} remain perfect across all six HEROS-profile runs. The results also show, however, that prompt design alone does not determine downstream explanation quality. The learned HEROS model that supplies the packet is itself part of the explanation system.

MUX11 provides the clearest positive case. Under stronger HEROS training, test accuracy rises from 81.2\% to 98.8\%, the selected model shrinks from 56 rules to 24 rules, the ideal-rule fraction rises from 18.8\% to 87.5\%, and mean matching rules drop from 4.86 to 2.24. Once the upstream model becomes cleaner and closer to the ideal MUX rule set, the downstream translation problem also becomes easier: \emph{Evidence Recall} rises from 94.4\% to 99.4\%, \emph{Rule Coverage} reaches 100.0\%, outputs become shorter, and \emph{Causal Overclaim Rate} drops by roughly half. MUX11 therefore shows that better HEROS training can directly improve explanation quality when it reduces packet density and removes redundant or non-ideal rules.

MUX20 shows the opposite case and keeps the paper from making too broad a claim. Stronger HEROS training improves HEROS test accuracy only modestly, from 60.1\% to 63.2\%, recovers no ideal rules, and leaves the selected model densely matched, with mean matching rules increasing slightly from 12.34 to 12.88. Under those conditions, \emph{Evidence Recall} falls from 73.7\% to 68.6\% and \emph{Rule Coverage} falls from 94.3\% to 93.0\%, even though \emph{Hallucination Rate} drops from 12.0\% to 8.7\% and both \emph{Audience Understandability} and \emph{Audience Technical Fit} improve slightly. The implication is that extra HEROS search budget alone is not enough. Downstream explanation quality improves only when better training also yields a simpler and more ideal rule model.

The experiment and audience comparisons remain informative once this upstream model-quality dimension is added. The minimal glossary in \emph{Experiment 2} behaves as a recurring multi-metric trade-off under both HEROS profiles. In four of the six dataset/profile comparisons it improves \emph{Evidence Recall}; in four it worsens \emph{Hallucination Rate}; in three it lowers \emph{Causal Overclaim Rate}; and in three it improves both \emph{Audience Understandability} and \emph{Audience Technical Fit}. At the same time, it consistently produces longer outputs and lower \emph{Flesch Reading Ease}, and it almost always raises \emph{Flesch--Kincaid Grade Level}. This pattern suggests that the glossary primarily increases phrasing latitude and explanatory detail rather than reliably selecting better evidence. The audience settings likewise form distinct error regimes rather than a quality ranking. \emph{Expert} prompting preserves the most packet evidence and receives the strongest technical-fit scores, but it is also the most prone to unsupported elaboration. \emph{Layman} prompting is the safest under the hallucination and causal-overclaim detectors, but it loses evidence coverage. \emph{Clinician} prompting stays between those extremes on coverage and fit, but its \emph{Causal Overclaim Rate} remains high relative to the other audiences. The readability results reinforce this interpretation by showing that the prompts do not yet produce a stable monotonic complexity ordering across audiences. We therefore treat the readability values as descriptive evidence about surface-form shifts rather than as validation that audience targeting is solved. Taken together, the findings do not support a single dominant configuration. They indicate instead that training profile, prompt experiment, and audience target each control different parts of the final behavior, and that evaluation has to consider all three dimensions jointly.

This study remains preliminary and has several limitations. First, it uses synthetic multiplexer benchmarks rather than biomedical datasets, so the present results should be read as a controlled proof of concept. Second, the six-run comparison still uses one seed and only two HEROS training profiles per dataset; it does not yet trace explanation behavior across a richer spectrum of model-selection checkpoints or repeated random initializations. Third, MUX11 and MUX20 were evaluated on stratified 50-instance subsets of their held-out folds rather than on the full held-out folds. Fourth, the judge metrics remain proxy signals rather than participant-based outcomes, and they should be interpreted with extra caution because an LLM from the same model family was used to judge LLM-generated explanations~\cite{liu2023geval,huang2025judge}. Fifth, the glossary intervention is itself subjective. We intentionally used minimal feature-only glossary entries because a more detailed glossary could leak the true multiplexer mechanism into the prompt, but that also means the study does not answer how specific a glossary should be. Finally, we do not report perturbation-based metrics such as comprehensiveness and sufficiency because computing them properly here requires rerunning HEROS under controlled feature-removal and keep-only perturbations~\cite{deyoung2020eraser}. Future work should extend the pipeline to biomedical rule learning problems, compare more HEROS training checkpoints and seeds, strengthen clinician-oriented non-causal prompting, add perturbation-based faithfulness analysis, and complement judge scoring with human evaluation using established explanation measures~\cite{hoffman2023measures,holzinger2020scs}.

\section{Code and Data Availability}
\begin{sloppypar}
The code branch corresponding to the work reported here is available at \url{https://github.com/UrbsLab/heros/tree/llm_exp}. The LLM explanation prototype used in this study is implemented under \path{evaluation/experiments/heros_llm}. The multiplexer datasets are available within the same repository under \path{evaluation/datasets/partitioned/multiplexer/}, including the predefined train/test split files for MUX6, MUX11, and MUX20 used in this study. The MUX6 experiments used the full 50-instance held-out test fold, while the MUX11 and MUX20 experiments used stratified 50-instance subsets sampled from their respective held-out test folds as described in the Methods section.
\end{sloppypar}

\section{Conclusion}
We presented a preliminary, script-first HEROS+LLM pipeline that treats the LLM as a translator of structured rule-based evidence rather than a generator of free-form reasons. Across six runs on MUX6, MUX11, and MUX20, the system maintained perfect \emph{Evidence Precision} and perfect \emph{Prediction--Explanation Agreement}, but the broader comparison showed that explanation quality depends jointly on upstream model quality, prompt design, and audience targeting. When stronger HEROS training simplified the selected rule model, as on MUX11, downstream evidence preservation improved sharply. When stronger training did not simplify the rule model, as on MUX20, explanation coverage did not improve even though classifier accuracy rose slightly. The glossary-style prompt experiment often improved audience-fit scores and sometimes lowered \emph{Causal Overclaim Rate}, but it also tended to lengthen text and could increase \emph{Hallucination Rate}. Audience customization exposed stable trade-offs rather than a single best output style. The main conclusion is therefore methodological as much as empirical: constrained LLM translation should be evaluated jointly with the learned rule model that supplies its packet, the prompt design that shapes the verbalization, and the audience for whom that verbalization is intended.

%%
%% The acknowledgments section is defined using the "acks" environment
%% (and NOT an unnumbered section). This ensures the proper
%% identification of the section in the article metadata, and the
%% consistent spelling of the heading.
\begin{acks}
The study was supported by Cedars Sinai Health Sciences University and NIH grants R01 AI173095, U01 AG066833, and P30 AG073105. Generative AI tools, including GPT-based systems, were used to assist with code scaffolding, analysis scripting, wording revisions, table preparation, and drafting portions of the manuscript text. All such material was reviewed, edited, and verified by the authors, who take full responsibility for the final content.

\end{acks}

\balance

%%
%% The next two lines define the bibliography style to be used, and
%% the bibliography file.
\bibliographystyle{ACM-Reference-Format}
\bibliography{references}

\end{document}
\endinput
%%
%% End of file `sample-sigconf.tex'.
